\section{Code Usage}
\label{sec:code}

\Bernardo{I will assume that \textsc{DSGRN\_utils} is merged with \textsc{DSGRN}. There are many utilities that I'm ignoring (equilibrium cells, etc.)}

Dynamic Signature of Gene Regulatory Networks (DSGRN) provides computational tools for analyzing network dynamics through the combinatorial decomposition of parameter space. A preliminary analysis allows the identification of multilevel models as in Figs~\ref{fig:multi1}(c) and~\ref{fig:multi2}(c) \cite{Cummins2017b}. These encompass all Boolean dynamics, and include the identification of multi-cell attractors and unstable strongly connected multi-cell sets in the discrete system. Further analysis provides a dynamical characterization of (sufficiently steep~\cite{rook_field:24}) continuous ODE models associated to the network structure.

In this section, we introduce readers to the expected output and interpretation of DSGRN, and include Jupyter notebook tutorials available with the software. DSGRN is installable via \texttt{
pip install DSGRN}. Our running example will be the toggle switch network from Section~\ref{subsec:TS}.

\subsection{Defining a Regulatory Network $\bRN$}

The first step is defining a regulatory network $\bRN=(V,E,\delta)$ (Definition~\ref{def:RN}) using DSGRN string-based specification format.  Each line in the network specification defines a node and its regulatory inputs associated via an algebraic expression. The format is
\begin{center}
  \texttt{node\_name : algebraic\_expression : optional\_parameter},
\end{center}
\begin{itemize}
  \item \texttt{node\_name}: Node identifier (e.g., gene or protein name)
  \item \texttt{algebraic\_expression}: Specifies network topology and edge signs. Activation of node $i$ is denoted by $i$ and repression by $\sim i$. The algebraic expression is given by a combination of products and sums, where $(i+j)(\sim k)$ indicates a node receives inputs from activators $i,j$ and repressor $k$.
  \item \texttt{optional\_modifiers}: Optional flags to restrict parameter space: \texttt{E} (Essential functions only), \texttt{B} (Boolean parameters only), \texttt{M} (Multi-Boolean), or combinations like \texttt{BE}.
\end{itemize}
\begin{rem}
  The algebraic expression serves two purposes: (1) it defines the network topology by identifying which edges exist and their signs (activation vs.\ repression), and (2) in the original DSGRN framework~\cite{Cummins2017b}, it determines the parameter space partition structure. In this work, we focus on Boolean parameters where the partition structure does not affect the parameter space.
\end{rem}
To define the network, we call the network constructor. For example,
\begin{minted}[fontsize=\small]{python}
import DSGRN
network_spec = """
v1 : v1 + v2
v2 : (~v1) + v2
"""
network = DSGRN.Network(network_spec)
DSGRN.DrawGraph(network) # plot the associated network
\end{minted}

\subsection{Parameter Graph}

The parameter graph $\PG = (\cP,\cE)$ (Section~\ref{sec:PG}) organizes all DSGRN parameters that are compatible with the regulatory network $\bRN$. The size of the parameter graph depends on the network structure and any specified modifiers. For our 2-node network:
\begin{minted}[fontsize=\small]{python}
parameter_graph = DSGRN.ParameterGraph(network)
print(parameter_graph.size())  # Output: 1600
\end{minted}

The parameter space can be restricted using modifiers in the network specification. For example, adding \texttt{B} to specify Boolean-only parameters reduces the parameter count, while \texttt{E} restricts to essential parameters.

\subsection{Morse Graph}
For a parameter $p \in \cP$, DSGRN computes the Morse graph $\MG(\cF)$ (Definition~\ref{def:DSGRN_MG}). We refer to \cite{rook_field:24} for the detailed algorithms that produce multivalued maps $\cF : \cX \rightrightarrows \cX$. To compute and plot the morse graph, and associated Morse sets (when possible), one can write
\begin{minted}[fontsize=\small]{python}
morse_graph, stg, graded_complex = DSGRN.ConleyMorseGraph(parameter)
DSGRN.PlotMorseGraph(morse_graph, label='a')
DSGRN.PlotMorseSets(morse_graph, stg, graded_complex)
\end{minted}
The \texttt{label} parameter controls how Morse graph nodes are annotated: \texttt{'c'} displays Conley indices (default), \texttt{'s'} shows the size (number of cells), and \texttt{'a'} provides annotations such as FP (fixed point), PO (periodic orbit), T (trivial), or M (more complex).

\subsection{A fully-connected 2-network}

Consider the following regulatory network $\bRN=(V,E,\delta)$ associated with the network Figure~\ref{fig:22network}(a)

\begin{minted}[fontsize=\small]{python}
net_spec = """
  v1 : v1 + v2
  v2 : (~v1)v2
"""
network = DSGRN.Network(net_spec)
parameter_graph = DSGRN.ParameterGraph(network)

print('Parameter graph size:', parameter_graph.size())
# output: Parameter graph size: 1600

par_index = 752
parameter = parameter_graph.parameter(par_index)

morse_graph, stg, graded_complex = DSGRN.ConleyMorseGraph(parameter)
DSGRN.PlotMorseGraph(morse_graph)
DSGRN.PlotMorseSets(morse_graph,stg,graded_complex)
\end{minted}

\subsection{Another example}

5-dimensional

\begin{minted}[fontsize=\small]{python}
net_spec = """
  v1 : v1 + (~v2) + (~v3)
  v2 : (~v1) + v2
  v3 : (~v2) + v3
  v4 : v2
  v5 : v3 + v4 + v5
"""

network = DSGRN.Network(net_spec)
parameter_graph = DSGRN.ParameterGraph(network)

print('Parameter graph size:', parameter_graph.size())
# output: Parameter graph size: 13608000000

par_index = 5103162287
parameter = parameter_graph.parameter(par_index)

morse_graph, stg, graded_complex = DSGRN.ConleyMorseGraph(parameter)
DSGRN.PlotMorseGraph(morse_graph)
\end{minted}